\subsection{Limitationen}\label{sec:limitationen}

Die erste und größte Grenze dieses Konzepts ist das Fehlen eines echten Nutzertests. Resteria beruht auf keinerlei Primärdaten aus der Zielgruppe; alle Annahmen zum Nutzungsverhalten, zur Attraktivität der Gamification-Mechanik und zur tatsächlichen Reduktion von Lebensmittelverschwendung bleiben deshalb hypothetisch. Diese Grenze betrifft nicht das gewählte Vorgehen, sondern die Geltung der Ergebnisse: Ob Resteria in der Praxis funktioniert, lässt sich erst nach einem realen Test beurteilen. Diese Grenze mildert dieser Bericht, indem zentrale Funktionen wie die Rettungspunkte aus überprüfter Verhaltensforschung stammen, etwa aus der Reminder-Wirkung, die \textcite{barkerWhatNudgeTechniques2021} belegen. Für die Umsetzung von Resteria ist deshalb ein Beta-Test vor dem Launch vorgesehen, wie ihn die Phasenplanung in~\autoref{sec:phasenplanung} bereits berücksichtigt.

Eine zweite Grenze betrifft das UI/UX-Mockup in Anhang B. Die Skizze ist KI-generiert und beruht nicht auf einer Rückmeldung realer Nutzer:innen, weshalb sich Aussagen zur tatsächlichen Bedienbarkeit von Resteria daraus nicht ableiten lassen. Diese Grenze mildert der Bericht, indem die Herkunft des Mockups transparent gekennzeichnet ist und sich die Darstellung auf die Kernfunktion Reste-Matching statt auf eine vollständige Ausarbeitung aller Bildschirme beschränkt. Vor einer realen Umsetzung sollte deshalb ein Usability-Test mit einem funktionsfähigen Prototyp folgen, der die Bedienführung tatsächlich prüft.

Die dritte Grenze liegt in der Wettbewerbsanalyse selbst. Sie untersucht mit Chefkoch.de, TikTok Food, Restegourmet, der \ac{BMEL}-App, leftovercooking und Foodsharing.de sechs Vertreter aus drei Plattform-Kategorien, aber nicht den gesamten Markt. Weitere, kleinere Reste-Matching-Apps oder regionale Rezept-Communitys bleiben unberücksichtigt, sodass sich die festgestellte Marktlücke nicht mit letzter Sicherheit verallgemeinern lässt. Der begrenzte Seitenumfang dieses Projektberichts rechtfertigt die gezielte statt erschöpfende Auswahl, bei der jede Kategorie mit ein bis drei Vertretern belegt ist. Vor einer realen Umsetzung von Resteria empfiehlt sich deshalb eine breitere Marktanalyse, die auch kleinere und internationale Anbieter einbezieht.
